\subsection{Radon-Nikodym Theorem}

\begin{theorem} \label{thm:radon-nikodym-sigma-finite} \hlt{Radon--Nikodym Theorem}\\
Let $\lambda,\mu$ be $\sigma$-finite measures on $(X,\mathcal{M})$ with $\lambda\ll\mu$.
\begin{enumerate}[label=(\roman*)]
\setlength{\itemsep}{0pt}
\item \hlt{Representation}: $\exists$ measurable $f:X\to[0,+\infty)$ such that $\lambda(E)=\int_Ef\,d\mu$ for $E\in\mathcal{M}$.
\item \hlt{Uniqueness}: Any two representing densities agree $\mu$-a.e.
\end{enumerate}
The $\mu$-a.e. class of $f$ is denoted by $\frac{d\lambda}{d\mu}$.
\begin{proof}
\begin{enumerate}[label=(\roman*)]
\setlength{\itemsep}{0pt}
\item Assume first $\lambda(X),\mu(X)<+\infty$. For $c>0$, let $(P_k,N_k)$ be Hahn for $\lambda-kc\mu$, set $A_k\defeq N_k\setminus\bigcup\limits_{j<k}A_j$, and set $B\defeq X\setminus\bigcup\limits_{k\geq1}A_k$. If $k\geq2$, then $A_k\subseteq N_k\cap P_{k-1}$.\\
Hence $E\subseteq A_k$ implies $(k-1)c\mu(E)\leq\lambda(E)\leq kc\mu(E)$, with the lower bound $0\leq\lambda(E)$ when $k=1$.\\
Since $B\subseteq P_k$ for $k\geq1$, $\lambda(B)\geq kc\mu(B)$ $\forall k\geq1$; finiteness and $\lambda\ll\mu$ give $\mu(B)=\lambda(B)=0$. Set $f_c\defeq\sum\limits_{k\geq1}(k-1)c\,\mathbf{1}_{A_k}$. Countable additivity over $(A_k)$ gives $0\leq\lambda(E)-\int_Ef_c\,d\mu\leq c\mu(E)$ for $E\in\mathcal{M}$.\\
Let $c_n\defeq2^{-n}$ and $f_n\defeq f_{c_n}$. For $E\in\mathcal{M}$, $\abs{\int_E(f_n-f_m)\,d\mu}\leq(c_n+c_m)\mu(E)$; by Proposition \ref{prop:l1-level-sets-null-and-indefinite-integrals}, $\abs{f_n-f_m}\leq c_n+c_m$ $\mu$-a.e. Thus $\sum\limits_n\norm{f_{n+1}-f_n}_1<+\infty$. By Tonelli, $f_n$ converges $\mu$-a.e. and in $L^1(\mu)$ to a measurable $f\geq0$. Letting $n\to\infty$ in the approximation inequality gives $\lambda(E)=\int_Ef\,d\mu$.\\
For the $\sigma$-finite case, choose a measurable partition $X=\bigsqcup\limits_rY_r$ with $\lambda(Y_r),\mu(Y_r)<+\infty$. Apply the finite case on each $Y_r$, obtaining $f_r$, and set $f\defeq\sum\limits_rf_r\mathbf{1}_{Y_r}$.\\
Then $\lambda(E)=\sum\limits_r\lambda(E\cap Y_r)=\int_Ef\,d\mu$; also $f<+\infty$ $\mu$-a.e. because $\int_{Y_r}f\,d\mu=\lambda(Y_r)<+\infty$ $\forall r$.
\item If $f,g$ represent $\lambda$, then $f,g\in L^1(\mu|_{Y_r})$ and $f=g$ $\mu$-a.e. on $Y_r$ by Proposition \ref{prop:l1-level-sets-null-and-indefinite-integrals}.\\
Taking the countable union of the exceptional sets gives $f=g$ $\mu$-a.e. on $X$.
\end{enumerate}
\hlt{Structural Commentary}: Hahn decompositions of $\lambda-c\mu$ locate density bands.\\
\hlt{Key Insight}: Summably refined band approximations converge without invoking later $L^1$ completeness.
\end{proof}
\end{theorem}

\begin{example} \label{ex:radon-nikodym-needs-sigma-finite-source} \hlt{A Non-$\sigma$-Finite Failure}\\
Let $X=[0,1]$, $\mathcal{M}=\mathcal{B}([0,1])$, $\mu=\#$, and $\lambda=m$. Since $\mu(E)=0\Leftrightarrow E=\emptyset$, $\lambda\ll\mu$.\\
If $\lambda(E)=\int_E f\,d\mu$ for some $f\geq0$, then $\forall x$, $0=\lambda(\{x\})=\int_{\{x\}}f\,d\mu=f(x)$, so $\lambda([0,1])=0$, contradiction.\\
Thus RN can fail when $\mu$ is not $\sigma$-finite.
\end{example}

\begin{example} \label{ex:countable-cocountable-radon-nikodym-failure} \hlt{Another Non-$\sigma$-Finite Failure}\\
Let $X$ be uncountable, $\mathcal{M}\defeq\{E\subseteq X:E\text{ or }X\setminus E\text{ countable}\}$, $\mu(E)\defeq\#E$ if $E$ is finite and $+\infty$ otherwise, and $\lambda(E)\defeq0$ if $E$ is countable and $+\infty$ otherwise. Then $\mu(E)=0\Leftrightarrow E=\emptyset$, so $\lambda\ll\mu$.\\
If $\lambda(E)=\int_Ef\,d\mu$ with $f\geq0$, then $\forall x$, $0=\lambda(\{x\})=\int_{\{x\}}f\,d\mu=f(x)$, forcing $\lambda(X)=0$, contradiction.\\
\hlt{Structural Commentary}: Singleton tests determine $f$, but not $\lambda(X)$ without $\sigma$-finiteness.\\
\hlt{Key Insight}: Non-$\sigma$-finite domination can lose global mass.
\end{example}

\begin{remark} \label{rem:rn-derivative-integrability-finiteness} \hlt{When the Radon--Nikodym Derivative Is Integrable}\\
In Theorem \ref{thm:radon-nikodym-sigma-finite}, $\frac{d\lambda}{d\mu}\in L^1(\mu)\Leftrightarrow\lambda(X)<+\infty$, since $\int_X\frac{d\lambda}{d\mu}\,d\mu=\lambda(X)$.
\end{remark}

\begin{theorem} \label{thm:radon-nikodym-extended-dominated-measure} \hlt{Extended-Valued Radon--Nikodym Theorem}\\
Let $\mu$ be $\sigma$-finite, let $\lambda$ be a measure on $(X,\mathcal{M})$, and assume $\lambda\ll\mu$.
\begin{enumerate}[label=(\roman*)]
\setlength{\itemsep}{0pt}
\item \hlt{Representation}: $\exists f:X\to[0,+\infty]$ measurable such that $\lambda(E)=\int_Ef\,d\mu$ for $E\in\mathcal{M}$.
\item \hlt{Uniqueness}: Any two representing densities agree $\mu$-a.e.
\end{enumerate}
\begin{proof}
\begin{enumerate}[label=(\roman*)]
\setlength{\itemsep}{0pt}
\item Assume first $\mu(X)<+\infty$ and set $\alpha\defeq\sup\{\mu(A):A\in\mathcal{M},\lambda(A)<+\infty\}$. Choose $(A_n)\subseteq\mathcal{M}$ with $\lambda(A_n)<+\infty$ and $\mu(A_n)\to\alpha$, and set $C_n\defeq\bigcup\limits_{j\leq n}A_j$, $N\defeq\bigcup\limits_nC_n$, and $P\defeq X\setminus N$. Then $\lambda|_N$ is $\sigma$-finite.\\
For $E\subseteq P$, $\lambda(E)<+\infty\Rightarrow\mu(E)=0$, since $\mu(E)>0$ would give $\mu(C_n\cup E)=\mu(C_n)+\mu(E)>\alpha$ eventually. Hence $\mu(E)>0\Rightarrow\lambda(E)=+\infty$, while $\mu(E)=0\Rightarrow\lambda(E)=0$ by $\lambda\ll\mu$.\\
By Theorem \ref{thm:radon-nikodym-sigma-finite}, $\exists g:N\to[0,+\infty)$ measurable such that $\lambda(E\cap N)=\int_{E\cap N}g\,d\mu$. Define $f=g$ on $N$ and $f=+\infty$ on $P$; preceding dichotomy gives $\lambda(E)=\int_Ef\,d\mu$. For $\sigma$-finite $\mu$, $X=\bigsqcup\limits_rY_r$, $\mu(Y_r)<+\infty$.\\
Let $f_r$ be the finite-$\mu$ density on $Y_r$ and set $f\defeq\sum\limits_rf_r\mathbf{1}_{Y_r}$. Then $\lambda(E)=\sum\limits_r\lambda(E\cap Y_r)=\int_Ef\,d\mu$.
\item If $f,h$ represent $\lambda$, fix $r,m,n\in\N$ and set $A_{m,n}^r\defeq Y_r\cap\{f\leq m\}\cap\{h\geq f+\frac{1}{n}\}$.\\
Since $\int_Af\,d\mu=\int_Ah\,d\mu$ for measurable $A\subseteq Y_r$ and $\int_{A_{m,n}^r}f\,d\mu<+\infty$, $\mu(A_{m,n}^r)=0$.\\
Taking the union over $r,m,n$ gives $\mu\{f<h\}=0$; symmetry gives $\mu\{h<f\}=0$.
\end{enumerate}
\hlt{Structural Commentary}: Maximal finite-$\lambda$ regions separate ordinary density from infinite density.\\
\hlt{Key Insight}: Outside the maximal $\sigma$-finite part, every positive-$\mu$ set has infinite $\lambda$-mass.
\end{proof}
\end{theorem}

\begin{theorem} \label{thm:lebesgue-decomposition-measures} \hlt{Lebesgue Decomposition Theorem}\\
Let $\lambda,\mu$ be $\sigma$-finite measures on $(X,\mathcal{M})$.
\begin{enumerate}[label=(\roman*)]
\setlength{\itemsep}{0pt}
\item $\exists$ measures $\lambda_s,\lambda_a$ with $\lambda=\lambda_s+\lambda_a$, $\lambda_s\perp\mu$, and $\lambda_a\ll\mu$.
\item If $(\tilde{\lambda}_s,\tilde{\lambda}_a)$ is another such pair, then $\tilde{\lambda}_s=\lambda_s$ and $\tilde{\lambda}_a=\lambda_a$.
\end{enumerate}
\begin{proof}
\begin{enumerate}[label=(\roman*)]
\setlength{\itemsep}{0pt}
\item Put $\rho\defeq\lambda+\mu$ and choose $f=\frac{d\lambda}{d\rho}$, $g=\frac{d\mu}{d\rho}$. With $A\defeq\{g=0\}$ and $B\defeq\{g>0\}$, set $\lambda_s(E)\defeq\lambda(E\cap A)$ and $\lambda_a(E)\defeq\lambda(E\cap B)$. Then $\lambda=\lambda_s+\lambda_a$. Since $\mu(A)=\int_Ag\,d\rho=0$, $\lambda_s\perp\mu$.\\
If $\mu(E)=0$, then $g=0$ $\rho$-a.e. on $E$, so $\rho(E\cap B)=0$ and $\lambda_a(E)=0$; hence $\lambda_a\ll\mu$.
\item Let $(\tilde{\lambda}_s,\tilde{\lambda}_a)$ be another such pair. Since $\lambda_s\perp\mu$ and $\tilde{\lambda}_a\ll\mu$, $\lambda_s\perp\tilde{\lambda}_a$; choose $X=C\cup D$ with $\lambda_s(C)=0$ and $\tilde{\lambda}_a(D)=0$. Then $\lambda_s(E)=\lambda_s(E\cap D)\leq\lambda(E\cap D)=\tilde{\lambda}_s(E\cap D)\leq\tilde{\lambda}_s(E)$.\\
By symmetry, $\lambda_s=\tilde{\lambda}_s$, hence $\lambda_a=\tilde{\lambda}_a$.
\end{enumerate}
\hlt{Structural Commentary}: Split $\lambda$ into $\mu$-visible and $\mu$-singular parts.\\
\hlt{Key Insight}: Use $\rho=\lambda+\mu$ and split where $\frac{d\mu}{d\rho}=0$.
\end{proof}
\end{theorem}

\begin{figure}[H]
  \centering
  \begin{tikzpicture}[
  scale=0.88,
  every node/.style={font=\scriptsize},
  box/.style={draw=black!50, line width=0.6pt, align=center, minimum width=1.62cm, minimum height=0.58cm, fill=white},
  formula/.style={box, minimum width=2.42cm, minimum height=0.70cm},
  posbox/.style={box, fill=blue!18},
  neg/.style={box, fill=red!10},
  ac/.style={box, fill=teal!12},
  sing/.style={box, fill=orange!16},
  arrow/.style={-{Latex[length=2mm]}, line width=0.55pt, black!60}
]
  \node[box] (charge) at (0,1.25) {charge\\$\lambda$};
  \node[posbox] (plus) at (-1.32,0.18) {$\lambda^+$\\positive};
  \node[neg] (minus) at (1.32,0.18) {$\lambda^-$\\negative};
  \node[formula] (jordan) at (0,-0.98) {$\lambda=\lambda^+-\lambda^-$\\$\abs{\lambda}=\lambda^++\lambda^-$};

  \draw[arrow] (charge) -- (plus);
  \draw[arrow] (charge) -- (minus);
  \draw[arrow] (plus) -- (jordan);
  \draw[arrow] (minus) -- (jordan);

  \node[box] (measure) at (4.60,1.25) {measure\\$\lambda$};
  \node[sing] (singular) at (3.28,0.18) {$\lambda_s\perp\mu$\\singular};
  \node[ac] (absolute) at (5.92,0.18) {$\lambda_a\ll\mu$\\visible};
  \node[formula] (leb) at (4.60,-0.98) {$\lambda=\lambda_s+\lambda_a$\\$\lambda_a(E)=\int_E f\,d\mu$};

  \draw[arrow] (measure) -- (singular);
  \draw[arrow] (measure) -- (absolute);
  \draw[arrow] (singular) -- (leb);
  \draw[arrow] (absolute) -- (leb);

  \draw[black!25, line width=0.5pt] (2.30,-1.48) -- (2.30,1.72);
\end{tikzpicture}

  \caption{Charges split by sign; measures split into $\lambda_s\perp\mu$ and $\lambda_a\ll\mu$ parts.}
  \label{fig:measure-decomposition-schematic}
\end{figure}

\begin{proposition} \label{prop:radon-nikodym-change-of-measure-calculus} \hlt{Radon--Nikodym Calculus}\\
Let all measures below be $\sigma$-finite.
\begin{enumerate}[label=(\roman*)]
\setlength{\itemsep}{0pt}
\item $\lambda\ll\mu$, $g\geq0$ measurable $\Rightarrow\int_Xg\,d\lambda=\int_Xg\frac{d\lambda}{d\mu}\,d\mu$.
\item $\nu\ll\lambda\ll\mu\Rightarrow\frac{d\nu}{d\mu}=\frac{d\nu}{d\lambda}\frac{d\lambda}{d\mu}$ $\mu$-a.e.
\item $\lambda_1,\lambda_2\ll\mu\Rightarrow\frac{d(\lambda_1+\lambda_2)}{d\mu}=\frac{d\lambda_1}{d\mu}+\frac{d\lambda_2}{d\mu}$ $\mu$-a.e.
\item $\lambda\ll\mu\ll\lambda\Rightarrow\frac{d\lambda}{d\mu}=\left(\frac{d\mu}{d\lambda}\right)^{-1}$ $\mu$-a.e.
\end{enumerate}
\begin{proof}
\begin{enumerate}[label=(\roman*)]
\setlength{\itemsep}{0pt}
\item For $g=\mathbf{1}_E$, $\int_Xg\,d\lambda=\lambda(E)=\int_E\frac{d\lambda}{d\mu}\,d\mu=\int_Xg\frac{d\lambda}{d\mu}\,d\mu$.\\
By linearity, identity holds for nonnegative simple $g$.\\
For measurable $g\geq0$, take simple $g_n\uparrow g$ and apply monotone convergence to both sides.
\item Let $h\defeq\frac{d\nu}{d\lambda}\frac{d\lambda}{d\mu}$. $\forall E\in\mathcal{M}$, (i) with $g=\mathbf{1}_E\frac{d\nu}{d\lambda}$ gives $\nu(E)=\int_E\frac{d\nu}{d\lambda}\,d\lambda=\int_Eh\,d\mu$.\\
Hence $h$ represents $\nu$ relative to $\mu$, so $h=\frac{d\nu}{d\mu}$ $\mu$-a.e.
\item Let $h\defeq\frac{d\lambda_1}{d\mu}+\frac{d\lambda_2}{d\mu}$. $\forall E\in\mathcal{M}$, $\int_Eh\,d\mu=\lambda_1(E)+\lambda_2(E)=(\lambda_1+\lambda_2)(E)$.\\
Hence $h$ represents $\lambda_1+\lambda_2$ relative to $\mu$, so uniqueness gives the identity.
\item Apply (ii) to $\mu\ll\lambda\ll\mu$. Since $\frac{d\mu}{d\mu}=1$, $1=\frac{d\mu}{d\lambda}\frac{d\lambda}{d\mu}$ $\mu$-a.e.\\
Since $\mu\ll\lambda$, $\frac{d\mu}{d\lambda}>0$ $\mu$-a.e.; divide to obtain $\frac{d\lambda}{d\mu}=\left(\frac{d\mu}{d\lambda}\right)^{-1}$ $\mu$-a.e.
\end{enumerate}
\hlt{Structural Commentary}: Radon--Nikodym derivatives obey chain and linearity rules.\\
\hlt{Key Insight}: Prove identities by integrating over arbitrary measurable sets, then invoke uniqueness.
\end{proof}
\end{proposition}

\begin{corollary} \label{cor:absolute-continuous-singular-zero-measure} \hlt{Absolutely Continuous and Singular Forces Zero}\\
Let $\lambda,\mu$ be measures and let $\eta$ be a charge.
\begin{enumerate}[label=(\roman*)]
\setlength{\itemsep}{0pt}
\item $\lambda\ll\mu$ and $\lambda\perp\mu$ imply $\lambda=0$.
\item $\abs{\eta}\perp\mu$ implies $\eta^+\perp\mu$.
\item $\abs{\eta}\perp\mu$ implies $\eta^-\perp\mu$.
\end{enumerate}
\begin{proof}
\begin{enumerate}[label=(\roman*)]
\setlength{\itemsep}{0pt}
\item Choose $X=A\cup B$ with $\lambda(A)=0$ and $\mu(B)=0$. Since $\lambda\ll\mu$, $\lambda(B)=0$, hence $\lambda(X)=0$.
\item Choose $X=A\cup B$ with $\abs{\eta}(A)=0$ and $\mu(B)=0$. Since $\eta^+\leq\abs{\eta}$, $\eta^+(A)=0$.
\item With the same $X=A\cup B$, $\eta^-(A)\leq\abs{\eta}(A)=0$ and $\mu(B)=0$; hence $\eta^-\perp\mu$.
\end{enumerate}
\hlt{Structural Commentary}: Absolute continuity and singular support are incompatible except at zero.\\
\hlt{Key Insight}: A measure supported on a null set must vanish when it also respects that null ideal.
\end{proof}
\end{corollary}

\begin{corollary} \label{cor:radon-nikodym-for-charges} \hlt{Radon--Nikodym for Charges}\\
Let $\mu$ be $\sigma$-finite and let $\lambda$ be a charge with $\lambda\ll\mu$. Then $\exists f\in\mathscr{L}^1(\mu)$ such that $\forall E\in\mathcal{M}$:
\begin{enumerate}[label=(\roman*)]
\setlength{\itemsep}{0pt}
\item $\lambda(E)=\int_E f\,d\mu$.
\item $\lambda^+(E)=\int_E f^+\,d\mu$.
\item $\lambda^-(E)=\int_E f^-\,d\mu$.
\item $\abs{\lambda}(E)=\int_E\abs{f}\,d\mu$.
\end{enumerate}
\begin{proof} Since $\abs{\lambda}\ll\mu$, $\lambda^\pm\ll\mu$. By Theorem \ref{thm:radon-nikodym-sigma-finite}, choose $f_\pm\geq0$ with $\lambda^\pm(E)=\int_Ef_\pm\,d\mu$.\\
Mutual singularity of $\lambda^+$ and $\lambda^-$ gives $f_+f_-=0$ $\mu$-a.e. Set $f\defeq f_+-f_-$.\\
Then $f^+=f_+$ and $f^-=f_-$ $\mu$-a.e., while $\int_X\abs f\,d\mu=\abs{\lambda}(X)<+\infty$.
\begin{enumerate}[label=(\roman*)]
\setlength{\itemsep}{0pt}
\item $\lambda(E)=\lambda^+(E)-\lambda^-(E)=\int_Ef\,d\mu$.
\item $\lambda^+(E)=\int_Ef_+\,d\mu=\int_Ef^+\,d\mu$.
\item $\lambda^-(E)=\int_Ef_-\,d\mu=\int_Ef^-\,d\mu$.
\item $\abs{\lambda}(E)=\lambda^+(E)+\lambda^-(E)=\int_E\abs f\,d\mu$.
\end{enumerate}
\hlt{Structural Commentary}: Signed absolute continuity is reduced to the two positive Jordan pieces.\\
\hlt{Key Insight}: Differentiate $\lambda^+$ and $\lambda^-$ separately, then subtract.
\end{proof}
\end{corollary}

\begin{example} \label{ex:stieltjes-basic-generators-rn-classification} \hlt{Radon--Nikodym Classification of Basic Stieltjes Generators}\\
Let $m$ be Lebesgue measure on $\mathcal{B}(\R)$, and let $\lambda_i$ be the Borel--Stieltjes measure generated by $g_i$.
\begin{enumerate}[label=(\roman*)]
\setlength{\itemsep}{0pt}
\item $g_1(x)=2x$: $\lambda_1=2m$, so $\lambda_1\ll m$, $\frac{d\lambda_1}{dm}=2$, $m\ll\lambda_1$, $\frac{dm}{d\lambda_1}=\frac{1}{2}$, and $\lambda_1(\R)=+\infty$.
\item $g_2(x)=\arctan x$: $\lambda_2(E)=\int_E(1+x^2)^{-1}\,dm$, so $\frac{d\lambda_2}{dm}=(1+x^2)^{-1}$, $\lambda_2(\R)=\pi$, $m\ll\lambda_2$, and $\frac{dm}{d\lambda_2}=1+x^2$.
\item $g_3=\mathbf{1}_{[0,+\infty)}$: $\lambda_3=\delta_0$, so $\lambda_3\perp m$, $\lambda_3(\R)=1$, $\lambda_3\not\ll m$, and $m\not\ll\lambda_3$.
\item $g_4(x)=x\mathbf{1}_{[0,+\infty)}(x)$: $\lambda_4(E)=m(E\cap[0,+\infty))$, so $\lambda_4\ll m$, $\frac{d\lambda_4}{dm}=\mathbf{1}_{[0,+\infty)}$, $m\not\ll\lambda_4$, and $\lambda_4(\R)=+\infty$.
\end{enumerate}
Thus $\lambda_1,\lambda_2,\lambda_4\ll m$, $\lambda_3\perp m$, only $\lambda_2,\lambda_3$ are finite, and $m\ll\lambda_i$ exactly for $i=1,2$.\\
\hlt{Structural Commentary}: Stieltjes generators split naturally into density, atom, and restriction examples.\\
\hlt{Key Insight}: A positive density makes null sets equivalent; an atom lives on a Lebesgue-null set.
\end{example}

\begin{theorem} \label{thm:charges-total-variation-banach} \hlt{Charges Form a Banach Space}\\
Let $\operatorname{ca}(X,\mathcal{M})\defeq\{\lambda:\mathcal{M}\to\R:\lambda\text{ is a charge}\}$ and $\norm{\lambda}_{\mathrm{var}}\defeq\abs{\lambda}(X)$.\\
With pointwise operations, $(\operatorname{ca}(X,\mathcal{M}),\norm{\cdot}_{\mathrm{var}})$ is Banach.
\begin{proof}
\begin{enumerate}[label=(\roman*)]
\setlength{\itemsep}{0pt}
\item Charge linearity gives closure under pointwise operations. The variation axioms follow from\\
$\abs{\lambda}(X)=0\Rightarrow\abs{\lambda}(E)=0$ $\forall E$, $\abs{a\lambda}=|a|\abs{\lambda}$, and $\sum\limits_j\abs{(\lambda+\eta)(A_j)}\leq\sum\limits_j\abs{\lambda(A_j)}+\sum\limits_j\abs{\eta(A_j)}$ over finite measurable partitions $(A_j)$.
\item Let $(\lambda_n)$ be variation-Cauchy.\\
Since $\abs{\lambda_n(E)-\lambda_m(E)}\leq\norm{\lambda_n-\lambda_m}_{\mathrm{var}}$, set $\lambda(E)\defeq\lim\limits_n\lambda_n(E)$; then $\lambda$ is finitely additive.\\
For fixed $n$ and finite disjoint $(A_j)$, $\sum\limits_j\abs{(\lambda-\lambda_n)(A_j)}\leq\limsup\limits_m\norm{\lambda_m-\lambda_n}_{\mathrm{var}}$.
\item For disjoint $(E_k)$, set $E\defeq\bigcup\limits_kE_k$, $R_K\defeq\bigcup\limits_{k>K}E_k$, and choose $n$ with $\limsup\limits_m\norm{\lambda_m-\lambda_n}_{\mathrm{var}}<\varepsilon$.\\
Since $\lambda_n(R_K)\to0$, for large $K$, $\abs{\lambda(R_K)}\leq\abs{\lambda_n(R_K)}+\abs{(\lambda-\lambda_n)(R_K)}<2\varepsilon$.\\
Hence $\lambda(R_K)\to0$, so $\lambda$ is countably additive.
\item Taking suprema in (ii) over finite partitions gives $\norm{\lambda-\lambda_n}_{\mathrm{var}}\leq\limsup\limits_m\norm{\lambda_m-\lambda_n}_{\mathrm{var}}\to0$.\\
Hence every variation-Cauchy sequence converges in variation to a charge.
\end{enumerate}
\hlt{Structural Commentary}: Total variation is the signed-measure analogue of an $L^1$ norm.\\
\hlt{Key Insight}: Variation-Cauchy controls every finite partition.
\end{proof}
\end{theorem}
